%! Author = lili
%! Date = 5/26/26

\section{What is RDF?}
\untit{Slides 1--4}
\Glspl{rdf} is intended for situations where information on the Web needs to be
processed by applications.
\defin{rdf}
It provides a common framework for exchanging information without loss of
meaning and enables \gls{linked-data} fpr publishing and interlinking data on the web.
% TODO loss of meaning?

The W3C (World Wide Web Consortium) standardizes the \gls{rdf}.

\section{Triple Model}
\untit{Slides 5--8}
\defin{triple-model}
\begin{itemize}
    \item \textbf{subject}: who or what we are talking about
    \item \textbf{predicate}: the relationship or property
    \item \textbf{object}: the thing the subject is related to
\end{itemize}

\begin{bsp}
todo
\end{bsp}
The same resource can appear in multiple triples, leading to a  NETWORK of
connected information which makes RDF powerful.
\begin{bsp}
    TODO Mona Lisa
\end{bsp}

\section{International Resource Identifiers (IRIs/URIs)}
\untit{Slides 8--9}
\defin{iri}

\begin{bsp} To \defiref{iri} \\
URLs are one form of IRI:
\begin{verbatim*}
    http://dbpedia.org/resource/Leonardo_da_Vinci
    http://xmlns.com/foaf/0.1/knows
    http://www.example.org/bob#me
\end{verbatim*}
\end{bsp}
\Glspl{iri} can appear in subject, predicate, or object position: \glspl{literal} are restricted
to object position.
\Glspl{iri} are also global identifiers, which means other people can reuse them. \Gls{rdf} is gnostic about what an \gls{iri} represents, only vocabularies provide meaning.
\begin{bsp} \textbf{Vocabularies} \\
FOAF, Dublin Core, or Schema.org
\end{bsp}
\section{Literals}
\untit{Slides 10--12}
\defin{literal}
\begin{bsp} To \defiref{literal} \\
strings, dates, numbers
\end{bsp}
\begin{notation}
    The $\^ \^$ symbol means ``has datatype''.
\end{notation}
\begin{bsp} \textbf{Literal with data type:}
    \begin{center}
        $"1990-07-04" \^ \^ <http://www.w3.org/2001/XMLSchema\#date>$
    \end{center}
    \textbf{Literal with language tag:}
    \begin{center}
        "Leonard de Vinci"@fr
    \end{center}
\end{bsp}

\begin{notte}
    \glspl{literal} may only appear in the object position if a triple.
\end{notte}


\begin{bsp} \textbf{Correct:}
\begin{center}
    :bob schema:birthDate "1990-07-04"^^xsd:date .
\end{center}
    Is interpreted as: Bob’s birth date is 1990-07-04 \\
\textbf{Incorrect:}
    \begin{center}
        "1990-07-04" schema:birthDate :bob
    \end{center}
    Is interpreted as: 1990-07-04 has birthDate Bob, because \glspl{literal} are values (like text, numbers, dates),
not entities that can act as subjects.
\end{bsp}


\section{Standard Format Turtle - Terse RDF Triple Language}
\untit{Slides 13--16}



\section{Standard Format JSON}
\untit{Slides 17--19}

\section{Flat Triplet Approach to RDF Graphs in JSON}
\untit{Slides 20--24}

\section{Comparison of Standard RDF Formats}
\untit{Slides 25--27}

\section{Summary}
\untit{Slides 28--29}

% Triple base database vs relational databases ??
% RDF dumps & sparql api
% owl mit sameAs wäre top für die iGEM Datenbank

